\documentclass[11pt]{article}

\usepackage[a4paper,margin=1in]{geometry}
\usepackage{amsmath,amssymb,amsfonts}
\usepackage{bm}
\usepackage{hyperref}
\usepackage{physics}
\usepackage{mathtools}
\usepackage{microtype}
\usepackage{enumitem}
\usepackage{graphicx}
\usepackage{courier}
\usepackage[T1]{fontenc}

\hypersetup{
  colorlinks=true,
  linkcolor=blue,
  citecolor=blue,
  urlcolor=blue
}

\title{\vspace{-1em}
Prime-Indexed Noncommutative Causal Dynamical Triangulations\\
as a Multiplicity-Theoretic Model of Quantum Spacetime}
\author{Multiplicity Foundation}
\date{\today}

\begin{document}
\maketitle

\begin{abstract}
We develop a Multiplicity-native synthesis of Noncommutative Geometry (NCG)
and Causal Dynamical Triangulations (CDT) by treating spacetime as a
prime-indexed, recursively updated multiplicity network of operator-valued
simplices. Building on previous drafts that treated NCG and CDT largely in
parallel, we introduce a coupled action density
\(
  s_i = S_{\mathrm{Regge},i} + S_{\mathrm{NCG},i}
       + \lambda \theta_p \varepsilon_i
\)
for simplex \(i\) in prime sector \(p\), and a dynamical feedback equation
for the multiplicity weights \(w_{p,i}(t)\) that drives the evolution of
the discrete geometry.\,[file:1] We implement this structure in a 2D CDT toy
simulator with four prime sectors \(p \in \{2,3,5,7\}\), providing a concrete
testbed for prime-structured quantum gravity and spectral-dimension
phenomenology.\,[file:1]
\end{abstract}
\newpage
\tableofcontents
\newpage
\section{Executive Summary}

\subsection{Motivation}

Multiplicity Theory views mathematical and physical objects as recursively
generated patterns of interactions indexed by primes, rather than as static
points on a manifold. In this ontology, quantum spacetime should emerge
from a prime-indexed network of relations, with identity and behavior
preserved through recursive feedback across scales.[file:1]

Noncommutative Geometry (NCG) suggests that at Planckian scales the
coordinates of spacetime become operators obeying non-trivial commutation
relations, while Causal Dynamical Triangulations (CDT) offers a discrete,
non-perturbative path integral over simplicial geometries with a causal
foliation.[web:13][web:16][web:19] Previous work integrated NCG and CDT into Multiplicity
largely as separate modules; here, we treat them as \emph{coupled
deformations} of a single prime-indexed multiplicity space.[file:1]

\subsection{Core Construction}

We define a \emph{prime-indexed noncommutative triangulated multiplicity
space} \(\mathbb{M}\) as the collection
\begin{equation}
  \mathbb{M} =
    \bigl\{ \bigl(w_{p,i}(t), \sigma_{p,i}, \hat{x}^\mu_{p,i} \bigr)
    \bigr\}_{p \in \mathcal{P},\; i \in I_p},
\end{equation}
where \(\mathcal{P}\) is the set of primes, \(\sigma_{p,i}\) encodes local
simplicial/combinatorial data, and \(\hat{x}^\mu_{p,i}\) are operator-valued
coordinates satisfying noncommutative relations.[file:1]
The dynamics is driven by:
\begin{enumerate}[label=(\alph*)]
  \item A Regge--NCG action density
  \(
    s_i = S_{\mathrm{Regge},i} + S_{\mathrm{NCG},i}
         + \lambda \theta_p \varepsilon_i
  \),
  where \(\varepsilon_i\) is the deficit angle, \(\theta_p\) the
  sector-dependent noncommutativity parameter, and \(\lambda\) a coupling.[file:1][web:13][web:20]
  \item A feedback equation for multiplicity weights \(w_{p,i}(t)\):
  \begin{equation}
    \frac{d w_{p,i}}{dt}
      = \alpha\bigl(\langle s \rangle_p - s_i\bigr)
      + \beta \sum_{p' \neq p} \theta_{pp'}
          \bigl(\langle s \rangle_{p'} - \langle s \rangle_p \bigr)
      - \gamma \bigl(w_{p,i} - w_0\bigr),
    \label{eq:feedback}
  \end{equation}
  where \(\langle s \rangle_p\) is the sector-averaged action density and
  \(\theta_{pp'}\) encodes cross-sector noncommutativity.[file:1]
\end{enumerate}

\subsection{Novelty}

Relative to standard NCG and CDT frameworks, the novelty is threefold:
\begin{itemize}
  \item NCG and CDT are not merely parallel modules but are coupled via
        the explicit term \(\lambda \theta_p \varepsilon_i\) in the local
        action density \(s_i\), making noncommutativity dynamically deform
        curvature rather than just modifying field interactions on a fixed
        background.[file:1][web:20]
  \item The degrees of freedom are prime-indexed simplices with
        recursively updated weights \(w_{p,i}(t)\), so prime structure enters
        directly into the geometry of spacetime and its evolution.[file:1]
  \item Tensor networks are reinterpreted as prime-labeled multiplicity
        networks, with phases and couplings determined by noncommutativity
        and Regge curvature, aligning with Multiplicity's broader program of
        prime-labeled recursion across mathematics and physics.[file:1]
\end{itemize}

\subsection{Implementation and Expected Signatures}

We implement the construction in a 2D CDT toy simulator with four prime
sectors \(p \in \{2,3,5,7\}\), constant sector-dependent noncommutativity
\(\theta_p\), and Euler integration of Eq.\,\eqref{eq:feedback} for the
weights \(w_{p,i}(t)\).[file:1] The simulator tracks:
\begin{enumerate}[label=(\roman*)]
  \item Sector-averaged weights \(\langle w \rangle_p(t)\),
  \item Sector-averaged curvature \(\langle \varepsilon \rangle_p(t)\),
  \item A proxy for the spectral dimension \(D_s(t)\) inferred from the
        curvature flow.[file:1][web:15][web:18]
\end{enumerate}
We expect:
\begin{itemize}
  \item Recovery of baseline 2D CDT spectral-dimension behavior when
        \(\beta = \lambda = 0\) and all \(\theta_p\) are equal, namely
        \(D_s \lesssim 2\) with smooth scale dependence.[web:15][web:18]
  \item Prime-sector competition, where sectors with favorable combinations
        of \(\theta_p\) and curvature \(\varepsilon_i\) dominate the late-time
        multiplicity distribution.
  \item Possible plateaus or oscillations in \(D_s\) at scales associated
        with dominant prime sectors, representing a characteristic
        ``Multiplicity signature'' absent in plain CDT.[file:1]
\end{itemize}

\section{Mathematical Framework}

\subsection{Noncommutative Spacetime Structure}

In Noncommutative Geometry, spacetime coordinates are promoted to
operators \(\hat{x}^\mu\) obeying
\begin{equation}
  [\hat{x}^\mu, \hat{x}^\nu]
    = i \theta^{\mu\nu},
\end{equation}
where \(\theta^{\mu\nu}\) is an antisymmetric matrix characterizing the
noncommutativity.\,[file:1][web:20] In the Multiplicity setting, we refine this
to \emph{prime- and simplex-resolved} operators:
\begin{equation}
  [\hat{x}^\mu_{p,i}, \hat{x}^\nu_{p',j}]
    = i \theta^{\mu\nu}_{pp',ij},
\end{equation}
where \(p,p'\) are primes and \(i,j\) label simplices within each
sector.[file:1]

In the simplest implementation considered here, we take
\(\theta^{\mu\nu}_{pp',ij}\) to factor as
\begin{equation}
  \theta^{\mu\nu}_{pp',ij} =
    \theta_0\,T^{\mu\nu}\,f_{pp'}\,g_{ij},
\end{equation}
where \(T^{\mu\nu}\) is a fixed antisymmetric tensor (e.g.\ a constant
spacelike plane), \(f_{pp'}\) encodes cross-prime coupling, and \(g_{ij}\)
encodes local structure (for the toy model we set \(g_{ij} \approx 1\)).[file:1]

\subsection{Moyal--Weyl Star Product and NCG Action}

Classically, fields \(\phi(x)\) on spacetime are replaced by fields
\(\hat{\phi}(\hat{x})\) on the noncommutative algebra, and multiplication is
deformed via the Moyal--Weyl star product
\begin{equation}
  (\phi \star \psi)(x)
    = \phi(x)\,
      \exp\!\left(
        \frac{i}{2} \theta^{\mu\nu}
        \overleftarrow{\partial_\mu}
        \overrightarrow{\partial_\nu}
      \right)\!
      \psi(x),
\end{equation}
so that the scalar action becomes
\begin{equation}
  S[\hat{\phi}]
    = \int d^4 x \sqrt{-g}\,
      \bigl(
        \hat{g}^{\mu\nu}\,
        \partial_\mu \hat{\phi} \star
        \partial_\nu \hat{\phi}
        - V(\hat{\phi})
      \bigr).\,[file:1][web:20]
\end{equation}
On a triangulated manifold, one must discretize both the geometry and the
star product. For the present 2D toy model we treat the NCG contribution
at each simplex \(i\) in sector \(p\) as an effective local term:
\begin{equation}
  S_{\mathrm{NCG},i}
    = \kappa\,\theta_p
    \quad\text{or more generally}\quad
  S_{\mathrm{NCG},i} = \kappa\,\theta_p^2\,G_i,
\end{equation}
where \(G_i\) is a discrete scalar-gradient energy built from neighboring
simplices, and \(\kappa\) is a coupling.[file:1] This already captures the key
intuition: larger \(\theta_p\) sectors carry a higher noncommutative
penalty unless compensated by geometric flow.

\subsection{CDT and the Regge Action}

Causal Dynamical Triangulations discretize spacetime as a piecewise linear
manifold built from simplices with a global time foliation. In Regge
calculus, curvature is concentrated at lower-dimensional ``hinges''
(triangle faces in 4D, vertices in 2D), and the Einstein--Hilbert action
is replaced by a sum over deficit angles \(\delta_j\) and volumes \(V_j\):[web:13][web:16][web:19]
\begin{equation}
  S_{\mathrm{Regge}} =
    \sum_{j \in T} \bigl( k\,\delta_j - \lambda_{\mathrm{cosm}}\,V_j \bigr),
\end{equation}
where \(k \equiv 1/(8\pi G)\) and \(\lambda_{\mathrm{cosm}}\) is the
cosmological constant.[web:13][web:16]

In 2D CDT, the triangulations have a global discrete time parameter and
spatial slices of fixed topology (e.g.\ circles). The spectral dimension
\(D_s\) of such geometries is known to be \(\leq 2\), with various studies
showing values around 2 for infinite triangulations and with matter
coupling.[web:15][web:18][web:21]

For simplex \(i\), we define a local Regge contribution
\begin{equation}
  S_{\mathrm{Regge},i} =
    \varepsilon_i\,A_i,
\end{equation}
where \(\varepsilon_i\) is the deficit angle and \(A_i\) the area of the
simplex.\,[file:1]

\subsection{Unified Local Action Density}

To couple NCG and CDT at the level of the local dynamics, we define at
each simplex \(i\) in prime sector \(p\) the local action density:
\begin{equation}
  s_i \equiv s_{p,i} =
    S_{\mathrm{Regge},i}
    + S_{\mathrm{NCG},i}
    + \lambda \theta_p \varepsilon_i,
  \label{eq:local-action}
\end{equation}
where:
\begin{itemize}
  \item \(S_{\mathrm{Regge},i} = \varepsilon_i A_i\),
  \item \(S_{\mathrm{NCG},i}\) is the effective NCG scalar contribution
        (constant or gradient-based),
  \item \(\lambda \theta_p \varepsilon_i\) is a new coupling that
        explicitly links noncommutativity (\(\theta_p\)) to curvature
        (\(\varepsilon_i\)).[file:1]
\end{itemize}
This term is precisely what closes the gap identified in the earlier
draft: NCG and CDT no longer live in separate channels but deform the same
object---the multiplicity-weighted prime-indexed triangulation.

\subsection{Sector Averages and Feedback}

For each prime sector \(p\), we define the multiplicity-weighted average
action
\begin{equation}
  \langle s \rangle_p(t) =
    \frac{1}{N_p}
    \sum_{i \in \mathrm{sector}\;p}
      w_{p,i}(t)\,s_{p,i}(t),
\end{equation}
where \(N_p\) is the number of simplices in sector \(p\). In practice, it
is natural to normalize by total weight rather than count:
\begin{equation}
  \langle s \rangle_p(t) =
    \frac{
      \sum_{i} w_{p,i}(t)\,s_{p,i}(t)
    }{
      \sum_{i} w_{p,i}(t)
    }.
\end{equation}
We also define sector-averaged curvature
\begin{equation}
  \langle \varepsilon \rangle_p(t) =
    \frac{
      \sum_{i} w_{p,i}(t)\,\varepsilon_{p,i}(t)
    }{
      \sum_{i} w_{p,i}(t)
    }.
\end{equation}

\subsection{Feedback Dynamics for Weights}

The core Multiplicity dynamics is encoded in the evolution equation for
the multiplicity weights \(w_{p,i}(t)\). We postulate the feedback equation:
\begin{equation}
  \frac{d w_{p,i}}{dt}
    = \alpha\bigl(\langle s \rangle_p - s_{p,i}\bigr)
    + \beta \sum_{p' \neq p}
        \theta_{pp'}
        \bigl(\langle s \rangle_{p'} - \langle s \rangle_p \bigr)
    - \gamma \bigl(w_{p,i} - w_0\bigr),
  \label{eq:feedback-full}
\end{equation}
where:
\begin{itemize}
  \item \(\alpha\) controls within-sector competition: simplices with lower
        local action \(s_{p,i}\) than the sector average gain weight.
  \item \(\beta\) controls cross-sector coupling, modulated by
        \(\theta_{pp'}\); differences in sector averages drive weight
        transfer among primes.
  \item \(\gamma\) controls relaxation toward the baseline multiplicity
        \(w_0\), ensuring numerical and conceptual stability.
\end{itemize}
This is the continuum form of the discrete update implemented in the
simulator.

\subsection{Linear Stability Near Equilibrium}

Write \(w_{p,i}(t) = w_0 + x_{p,i}(t)\) and consider small deviations
\(x_{p,i} \ll w_0\). If the sector averages vary slowly relative to
\(x_{p,i}\), the dominant local term is the decay component:
\begin{equation}
  \frac{dx_{p,i}}{dt} \approx -\gamma x_{p,i}.
\end{equation}
Using explicit Euler with time step \(dt\), we have
\begin{equation}
  x^{(n+1)}_{p,i} =
    x^{(n)}_{p,i} + dt\,\frac{dx_{p,i}}{dt}
    \approx (1 - \gamma dt)\,x^{(n)}_{p,i}.
\end{equation}
Stability of this mode requires \(|1 - \gamma dt| < 1\), which gives the
condition
\begin{equation}
  0 < \gamma\,dt < 2.
\end{equation}
For the simulator choice \(dt = 0.05\), this permits \(\gamma\) up to
approximately 40, while the current experiments use \(\gamma \lesssim 0.5\),
comfortably inside the stable region.[code_file:1]

The full coupled system introduces additional eigenmodes from the
cross-sector and within-sector terms, but as long as the effective
coefficients \(\alpha, \beta, \lambda,\theta_0\) remain \(O(1)\), and the
weights are clamped to \(w_{p,i} \geq \epsilon\), the explicit Euler scheme
is robust in the toy model regime.

\section{Prime-Indexed Tensor Networks}

Tensor networks in Multiplicity represent quantum states and entanglement
across spacetime. In the NCG+CDT+Multiplicity setting, we refine the
state:
\begin{equation}
  \hat{\Psi}(t) =
    \sum_{(p,i),(p',j)}
      \hat{T}_{(p,i),(p',j)}(t)\,
      \hat{\Psi}_{p,i} \otimes \hat{\Psi}_{p',j}\,
      e^{i(\theta_{p,i}(t) + \theta_{p',j}(t))},
\end{equation}
where:
\begin{itemize}
  \item \(\hat{T}_{(p,i),(p',j)}\) encodes couplings between simplices in
        prime sectors \(p\) and \(p'\),
  \item \(\theta_{p,i}(t)\) are phase factors influenced by the
        noncommutative structure,
  \item the indices themselves are prime-resolved, consistent with the
        multiplicity ontology.[file:1]
\end{itemize}
CDT constrains which pairs \((p,i),(p',j)\) can couple (respecting
causality and adjacency in the simplicial complex), while NCG determines
phases and local star-product corrections. The multiplicity weights
\(w_{p,i}(t)\) can be interpreted as modulating the amplitude or norm
associated with each local tensor, naturally linking the geometric and
quantum structures.

\section{2D CDT Toy Model Implementation}

\subsection{Setup}

The current implementation focuses on a 2D CDT-inspired toy model with
four prime sectors \(p \in \{2,3,5,7\}\), each with \(N_p\) simplices and
sector-dependent noncommutativity \(\theta_p = \theta_0 \times
\{1.0,1.5,0.8,2.0\}\).\,[file:1] Each simplex in sector \(p\) has:
\begin{itemize}
  \item A weight \(w_{p,i}(t)\),
  \item A deficit angle \(\varepsilon_{p,i}\),
  \item An area \(A_{p,i}\),
  \item A position \((x_{p,i},y_{p,i})\) in a causal-strip layout,
  \item A local NCG contribution \(S_{\mathrm{NCG},i}\).
\end{itemize}

The local action density is implemented as
\begin{equation}
  s_{p,i}(t)
    = \varepsilon_{p,i}(t)\,A_{p,i}
      + S_{\mathrm{NCG},i}(t)
      + \lambda \theta_p \varepsilon_{p,i}(t),
\end{equation}
and the feedback equation \eqref{eq:feedback-full} is advanced with an
explicit Euler step of size \(dt = 0.05\).[file:1]

\subsection{Representative Code Snippet}

The following JavaScript snippet, extracted and simplified from the
simulator, illustrates the implementation of the local action density and
feedback update (comments removed for brevity):

\begin{verbatim}
function actionDensity(s, theta_p, lambda) {
  const regge   = s.epsilon * s.area;
  const ncg     = s.ncg;              // effective NCG term
  const coupling = lambda * theta_p * s.epsilon;
  return regge + ncg + coupling;
}

function sectorAvg(sector, lambda) {
  const tp = sector.theta;
  let sum = 0, wsum = 0;
  for (const s of sector.simplices) {
    sum  += s.w * actionDensity(s, tp, lambda);
    wsum += s.w;
  }
  return wsum > 0 ? sum / wsum : 0;
}

function stepSim() {
  const alpha  = parseFloat(slAlpha.value);
  const beta   = parseFloat(slBeta.value);
  const gamma  = parseFloat(slGamma.value);
  const lambda = parseFloat(slLambda.value);
  const w0     = parseFloat(slW0.value);
  const theta0 = parseFloat(slTheta.value);
  const dt     = 0.05;

  simState.sectors.forEach((sec, pi) => {
    sec.theta = theta0 * THETA_SCALE[pi];
  });

  const avgs = simState.sectors.map(sec => sectorAvg(sec, lambda));

  simState.sectors.forEach((sec, pi) => {
    for (const s of sec.simplices) {
      const si = actionDensity(s, sec.theta, lambda);
      const T1 = alpha * (avgs[pi] - si);
      let T2 = 0;
      simState.sectors.forEach((sec2, pj) => {
        if (pj !== pi) {
          const theta_pp = theta0 * THETA_SCALE[pi] * THETA_SCALE[pj];
          T2 += beta * theta_pp * (avgs[pj] - avgs[pi]);
        }
      });
      const T3 = -gamma * (s.w - w0);
      s.w = Math.max(0.01, s.w + dt * (T1 + T2 + T3));
    }
  });
}
\end{verbatim}

This realizes the feedback equation with:
\begin{itemize}
  \item \(S_{\mathrm{Regge},i} = \varepsilon_i A_i\),
  \item \(S_{\mathrm{NCG},i} \equiv s.\mathrm{ncg}\),
  \item \(\lambda \theta_p \varepsilon_i\) as an explicit coupling,
  \item \(\theta_{pp'} \propto \theta_0 \theta_p \theta_{p'}\).
\end{itemize}

\subsection{Baseline and Multiplicity Effects}

To validate the model, one considers the following regimes:
\begin{enumerate}[label=(\alph*)]
  \item \textbf{Baseline CDT-like regime}: set \(\beta = 0\), \(\lambda = 0\),
        small \(\gamma\), and all \(\theta_p\) equal. The system should
        approximate plain 2D CDT plus a trivial scalar field, with spectral
        dimension \(D_s \lesssim 2\) and smooth scale dependence.[web:15][web:18]
  \item \textbf{Multiplicity regime}: restore \(\lambda > 0\) and
        sector-dependent \(\theta_p\), then vary \(\beta\). Here one can
        observe:
        \begin{itemize}
          \item Which prime sector dominates the late-time average weight
                \(\langle w \rangle_p(t)\),
          \item How the curvature \(\langle \varepsilon \rangle_p(t)\) flows
                for each sector,
          \item Whether \(D_s(t)\) exhibits plateaus or oscillations at
                scales associated with dominant sectors.
        \end{itemize}
\end{enumerate}

\section{Outlook and Validation Path}

\subsection{Prime Indexing and Dynamic Reassignment}

In the current prototype, sector labels are assigned by hand. A more
fundamental Multiplicity implementation would derive the prime sector of
each simplex from a global integer label via prime factorization, allowing
dynamic reassignment of degrees of freedom across sectors as the
triangulation evolves.[file:1]

\subsection{Phase Diagrams and Spectral Signatures}

By scanning \((\lambda,\beta)\) space and recording:
\begin{itemize}
  \item The fraction of late-time steps where each prime sector is
        dominant,
  \item The late-time mean and variance of \(D_s\),
\end{itemize}
one can map out phase diagrams indicating regimes of sector dominance,
coexistence, and oscillatory behavior. Comparing these diagrams to the
\(\beta=\lambda=0\) baseline quantifies the ``Multiplicity effect'' on
spectral dimension and curvature flow.[file:1]

\subsection{Higher Dimensions and Star-Product Refinement}

Future work includes:
\begin{itemize}
  \item Extending to 2+1 dimensions, using the full Regge action with
        2-simplices as hinges and 3-simplices as building blocks.[web:16][web:22]
  \item Implementing a more faithful discretization of the Moyal star
        product on the triangulation, possibly via dual graphs or matrix
        models of noncommutative field theory.[web:14][web:17][web:20]
  \item Connecting the weight evolution \(w_{p,i}(t)\) to PhaseMirror-style
        entanglement spectra across primes, using Schmidt decompositions of
        the prime-indexed tensor network.\,[file:1]
\end{itemize}

\subsection{Conceptual Significance}

The construction presented here realizes, in a concrete and simulatable
setting, the idea that spacetime is not a static manifold but a
prime-indexed multiplicity space whose geometry and noncommutativity are
jointly determined by recursive feedback dynamics. NCG and CDT become
different facets of a single multiplicity-theoretic object, rather than
external modules grafted onto GR. This paves the way for a genuinely
arithmetic, recursively stable approach to quantum gravity and cosmology
within the Multiplicity program.[file:1][web:15][web:18]

\appendix
\section*{Mathematical Appendix}
\addcontentsline{toc}{section}{Mathematical Appendix}

\section{Operator-Theoretic Setting and Notation}

\subsection{Prime-Indexed Hilbert Decomposition}

Let \(\mathcal{P}\) denote the set of prime numbers and let
\(\{I_p\}_{p \in \mathcal{P}}\) be index sets labelling the simplices in
each prime sector. For each \(p \in \mathcal{P}\), we consider a
finite-dimensional Hilbert space
\[
  \mathcal{H}_p
  := \ell^2(I_p)
  = \left\{
      \psi = (\psi_i)_{i \in I_p}
      : \sum_{i \in I_p} |\psi_i|^2 < \infty
    \right\},
\]
equipped with the standard inner product
\(
  \langle \psi, \varphi \rangle_p
  = \sum_{i \in I_p} \overline{\psi_i} \varphi_i.
\)[file:1]

The global multiplicity Hilbert space is the orthogonal direct sum
\begin{equation}
  \mathcal{H}
  := \bigoplus_{p \in \mathcal{P}} \mathcal{H}_p,
\end{equation}
so that each \(\Psi \in \mathcal{H}\) decomposes as
\(
  \Psi = \bigoplus_{p \in \mathcal{P}} \Psi^{(p)}
\)
with \(\Psi^{(p)} \in \mathcal{H}_p\).[file:1]

For each sector \(p\), we define:
\begin{itemize}
  \item The \emph{weight vector}
  \(
    w^{(p)}(t) = \{w_{p,i}(t)\}_{i \in I_p} \in (0,\infty)^{I_p}
  \),
  \item The \emph{local action density} vector
  \(
    s^{(p)}(t) = \{ s_{p,i}(t) \}_{i \in I_p},
  \)
  where
  \begin{equation}
    s_{p,i}(t)
      = S_{\mathrm{Regge},p,i}(t)
        + S_{\mathrm{NCG},p,i}(t)
        + \lambda\,\theta_p\,\varepsilon_{p,i}(t).
  \end{equation}[file:1]
\end{itemize}

We also define the diagonal multiplication operators
\(
  W_p(t), S_p(t) : \mathcal{H}_p \to \mathcal{H}_p
\)
by
\begin{equation}
  (W_p(t)\psi)_i = w_{p,i}(t)\psi_i,
  \qquad
  (S_p(t)\psi)_i = s_{p,i}(t)\psi_i.
\end{equation}
Their direct sums yield bounded operators on \(\mathcal{H}\):
\(
  W(t) = \bigoplus_p W_p(t),
  \ S(t) = \bigoplus_p S_p(t).
\)

\subsection{Sector-Average Operators}

For each sector, we define the scalar
\begin{equation}
  \langle s \rangle_p(t)
    := \frac{\sum_{i \in I_p} w_{p,i}(t)\,s_{p,i}(t)}
             {\sum_{i \in I_p} w_{p,i}(t)}
  = \frac{
      \langle w^{(p)}(t), s^{(p)}(t) \rangle_p
    }{
      \|w^{(p)}(t)\|_1
    },
\end{equation}
where \(\|\cdot\|_1\) is the \(\ell^1\)-norm on \(\mathbb{R}^{I_p}\).[file:1]
We also introduce the rank-one projection operators
\begin{equation}
  P_p(t) : \mathcal{H}_p \to \mathcal{H}_p,
  \qquad
  P_p(t)\psi
    = \frac{\langle w^{(p)}(t), \psi \rangle_p}{\|w^{(p)}(t)\|_2^2}
      \, w^{(p)}(t),
\end{equation}
with \(\|\cdot\|_2\) the \(\ell^2\)-norm. These encode the ``direction''
selected by the current multiplicity in each prime sector.

\section{Feedback Dynamics as an Operator Equation}

\subsection{Componentwise Form and Matrix Representation}

Recall the feedback equation for the weights:
\begin{equation}\label{eq:feedback-appendix}
  \frac{d w_{p,i}}{dt}
    = \alpha\bigl(\langle s \rangle_p - s_{p,i}\bigr)
    + \beta \sum_{p' \neq p} \theta_{pp'}
        \bigl(\langle s \rangle_{p'} - \langle s \rangle_p\bigr)
    - \gamma \bigl(w_{p,i} - w_0\bigr),
\end{equation}
with parameters \(\alpha,\beta,\gamma,w_0,\lambda,\theta_{pp'}\) as in the
main text.[file:1]

Fix \(t\) and view \eqref{eq:feedback-appendix} as an ODE in the vector
\(w(t) \in \mathbb{R}^{\bigsqcup_{p} I_p}\). We may write
\begin{equation}
  \frac{d}{dt} w_{p,i}
    = -\gamma w_{p,i}
      + \gamma w_0
      - \alpha s_{p,i}
      + \alpha \langle s \rangle_p
      + \beta \sum_{p' \neq p} \theta_{pp'}
          \bigl(\langle s \rangle_{p'} - \langle s \rangle_p \bigr).
\end{equation}
The terms involving \(\langle s \rangle_p\) are sector-constant and can be
viewed as collective driving terms. For fixed \(s_{p,i}\) and \(\langle s
\rangle_p(t)\), the evolution is affine-linear in \(w_{p,i}\).

At the operator level, one can write on each sector:
\begin{equation}
  \frac{d}{dt} W_p(t)
    \approx -\gamma \bigl( W_p(t) - w_0\,\mathbf{1}_p \bigr)
      - \alpha S_p(t)
      + \alpha \langle s \rangle_p(t)\,\mathbf{1}_p
      + \mathcal{C}_p(t),
\end{equation}
where \(\mathbf{1}_p\) is the identity on \(\mathcal{H}_p\), and
\(\mathcal{C}_p(t)\) is a scalar multiple of \(\mathbf{1}_p\) capturing
the cross-sector coupling:
\begin{equation}
  \mathcal{C}_p(t)
    = \beta \sum_{p' \neq p}
        \theta_{pp'}
        \bigl(\langle s \rangle_{p'}(t) - \langle s \rangle_p(t) \bigr)
        \,\mathbf{1}_p.
\end{equation}

\subsection{Linearization Around Equilibrium}

Let \(w_{p,i}(t) = w_0 + x_{p,i}(t)\) and assume a static background
\(\bar{s}_{p,i}\) with \(\langle \bar{s} \rangle_p\) independent of time.
Then the linearization of \eqref{eq:feedback-appendix} around the fixed
point \(x_{p,i}=0\) yields:
\begin{equation}
  \frac{dx_{p,i}}{dt}
    = -\gamma x_{p,i}
      + \alpha\bigl(\langle s \rangle_p - \bar{s}_{p,i} \bigr)
      + \beta \sum_{p' \neq p} \theta_{pp'}
          \bigl(\langle s \rangle_{p'} - \langle s \rangle_p \bigr)
      + \mathcal{O}(\|x\|^2),
\end{equation}
where \(\langle s \rangle_p = \langle \bar{s} \rangle_p + \mathcal{O}(x)\).
Expanding to first order in \(x\), we obtain a system of the form:
\begin{equation}
  \frac{dx}{dt}
    = -\gamma x + L x + b,
\end{equation}
where:
\begin{itemize}
  \item \(x \in \mathbb{R}^{\bigsqcup_p I_p}\) collects all deviations,
  \item \(L\) is a matrix capturing the dependence of the averages
        \(\langle s \rangle_p\) on \(x\),
  \item \(b\) is a constant vector.
\end{itemize}
The leading-order stability is dictated by the eigenvalues of
\(-\gamma I + L\).

\begin{lemma}[Scalar mode stability]
If we ignore the dependence of \(\langle s \rangle_p\) on \(x\) (i.e. treat
them as constants), then for each component \(x_{p,i}\) the linearized
equation reduces to
\begin{equation}
  \frac{dx_{p,i}}{dt} = -\gamma x_{p,i},
\end{equation}
which has solution \(x_{p,i}(t) = x_{p,i}(0)e^{-\gamma t}\). In
particular, the equilibrium \(x_{p,i}=0\) is exponentially stable with rate
\(\gamma\).
\end{lemma}

\begin{proof}
Immediate integration of the scalar ODE
\(
  \dot{x} = -\gamma x
\)
yields
\(
  x(t) = e^{-\gamma t} x(0)
\),
so \(|x(t)| \le e^{-\gamma t} |x(0)|\).
\end{proof}

\subsection{Explicit Euler and Operator Norm Bound}

Consider the explicit Euler discretization with step \(dt>0\):
\begin{equation}
  x^{(n+1)}_{p,i}
    = x^{(n)}_{p,i}
      + dt\,\frac{dx_{p,i}}{dt}\bigg|_{t = n dt}.
\end{equation}
In the scalar approximation \(\dot{x} = -\gamma x\), we have
\begin{equation}
  x^{(n+1)} = (1 - \gamma dt)\,x^{(n)}.
\end{equation}

\begin{lemma}[Scalar Euler stability]
The explicit Euler scheme for \(\dot{x} = -\gamma x\) is stable in the
sense \(|x^{(n)}| \to 0\) as \(n\to\infty\) if and only if
\begin{equation}
  |1 - \gamma dt| < 1,
  \quad\text{i.e.}\quad
  0 < \gamma dt < 2.
\end{equation}
\end{lemma}

\begin{proof}
We have
\(
  x^{(n)} = (1 - \gamma dt)^n x^{(0)}.
\)
This converges to \(0\) as \(n\to\infty\) precisely when
\(|1 - \gamma dt|<1\), which is equivalent to
\(
  -1 < 1 - \gamma dt < 1
  \iff 0 < \gamma dt < 2.
\)
\end{proof}

\begin{proposition}[Operator norm bound for Euler step]
Let \(\mathcal{T} : \mathcal{H} \to \mathcal{H}\) be the linear operator
defined by
\(
  \mathcal{T} = I - \gamma dt\,I + dt\,L
\),
where \(L\) is the linearization matrix of the feedback dynamics. Suppose
\(\|L\|_{\mathcal{B}(\mathcal{H})} \leq C_L\) and choose \(dt>0\) such that
\begin{equation}
  \gamma dt > C_L dt,
  \qquad
  \gamma dt < 2 - C_L dt.
\end{equation}
Then
\begin{equation}
  \|\mathcal{T}\|_{\mathcal{B}(\mathcal{H})}
    \leq 1 - (\gamma - C_L) dt < 1,
\end{equation}
and the discrete-time evolution is a strict contraction in operator norm.
\end{proposition}

\begin{proof}
We have
\(
  \mathcal{T} = I + dt(-\gamma I + L)
\),
so
\[
  \|\mathcal{T}\|
    \leq \|I\| + dt\|-\gamma I + L\|
    \leq 1 + dt(\gamma + C_L).
\]
This upper bound is not sharp. A refined argument uses the spectral radius
\(\rho(L)\) and the fact that the eigenvalues of \(\mathcal{T}\) are
\(1 + dt(-\gamma + \lambda_L)\) where \(\lambda_L\) ranges over the
spectrum of \(L\). If all eigenvalues satisfy
\(|1 + dt(-\gamma + \lambda_L)| < 1\), then
\(\|\mathcal{T}\| < 1\). One sufficient condition is
\(|-\gamma + \lambda_L| < \frac{2}{dt}\) for all \(|\lambda_L| \le C_L\),
which is equivalent to the inequalities in the statement. Under these
conditions, \(\mathcal{T}\) is a strict contraction in operator norm.
\end{proof}

In the simulator regime, \(dt = 0.05\) and \(\gamma \lesssim 0.5\), while
\(\alpha,\beta,\lambda,\theta_0\) are \(O(1)\), so the linearization
constant \(C_L\) is modest and the discretization is numerically stable
for the exploratory runs considered.[code_file:1][file:1]

\section{Bounds on Local Action Operators}

\subsection{Uniform Bounds on Regge and NCG Contributions}

Assume that in the 2D toy model the deficit angles and areas satisfy
\begin{equation}
  0 < \varepsilon_{\min} \leq \varepsilon_{p,i}(t) \leq \varepsilon_{\max},
  \qquad
  0 < A_{\min} \leq A_{p,i}(t) \leq A_{\max},
\end{equation}
and that the effective NCG term satisfies
\begin{equation}
  |S_{\mathrm{NCG},p,i}(t)| \leq K_{\mathrm{NCG}},
\end{equation}
for all \(p,i,t\). Furthermore, let \(|\theta_p| \leq \Theta_{\max}\) and
\(|\lambda| \leq \Lambda_{\max}\). Then the local action density satisfies
the uniform bound
\begin{equation}
  |s_{p,i}(t)|
    \leq \varepsilon_{\max} A_{\max}
         + K_{\mathrm{NCG}}
         + \Lambda_{\max} \Theta_{\max} \varepsilon_{\max}
    =: K_s.
\end{equation}

\begin{lemma}[Boundedness of \(S_p(t)\)]
Under the above assumptions, the diagonal operator \(S_p(t)\) on
\(\mathcal{H}_p\) is bounded with operator norm
\(
  \|S_p(t)\|_{\mathcal{B}(\mathcal{H}_p)} \leq K_s
\),
uniformly in \(t\).
\end{lemma}

\begin{proof}
For any \(\psi \in \mathcal{H}_p\),
\begin{align*}
  \|S_p(t)\psi\|_2^2
    &= \sum_{i \in I_p} |s_{p,i}(t)|^2 |\psi_i|^2
     \leq K_s^2 \sum_{i \in I_p} |\psi_i|^2
     = K_s^2 \|\psi\|_2^2.
\end{align*}
Thus \(\|S_p(t)\| \leq K_s\).
\end{proof}

Consequently, the global operator \(S(t) = \bigoplus_p S_p(t)\) is bounded
with \(\|S(t)\| \leq K_s\).

\subsection{Bounds on Sector-Average Functionals}

Define the linear functionals
\begin{equation}
  \mathcal{L}_p(t) : \mathcal{H}_p \to \mathbb{R},
  \qquad
  \mathcal{L}_p(t)(\psi)
    := \frac{\langle w^{(p)}(t), \psi \rangle_p}
             {\|w^{(p)}(t)\|_1},
\end{equation}
so that \(\langle s \rangle_p(t) = \mathcal{L}_p(t)\bigl(s^{(p)}(t)\bigr)\).
Assume there exist constants \(0 < m_w \leq M_w < \infty\) such that
\[
  m_w \leq w_{p,i}(t) \leq M_w
  \quad\text{for all } p,i,t.
\]
Then \(\|w^{(p)}(t)\|_1 \geq m_w |I_p|\) and
\(\|w^{(p)}(t)\|_2 \leq M_w \sqrt{|I_p|}\). By Cauchy--Schwarz,
\begin{align*}
  |\mathcal{L}_p(t)(\psi)|
    &= \frac{
         |\langle w^{(p)}(t), \psi \rangle_p|
       }{
         \|w^{(p)}(t)\|_1
       }
     \leq \frac{
         \|w^{(p)}(t)\|_2 \|\psi\|_2
       }{
         \|w^{(p)}(t)\|_1
       }
     \leq \frac{
         M_w \sqrt{|I_p|}
       }{
         m_w |I_p|
       } \|\psi\|_2
     = \frac{M_w}{m_w} |I_p|^{-1/2} \|\psi\|_2.
\end{align*}
Thus
\begin{equation}
  \|\mathcal{L}_p(t)\|
    \leq \frac{M_w}{m_w} |I_p|^{-1/2}.
\end{equation}

In particular, for fixed finite \(I_p\), the sector-average maps are
uniformly bounded linear functionals.

\section{Spectral Dimension Proxy and Bounds}

In the toy model, the spectral dimension \(D_s(t)\) is not computed via a
full diffusion process but via a curvature-based proxy:
\begin{equation}
  D_s(t)
    \approx 2.0 - c\,\bar{\varepsilon}(t),
\end{equation}
where \(\bar{\varepsilon}(t)\) is a global curvature average and \(c>0\)
is chosen such that \(D_s\) remains in the interval \([1,2]\).[file:1][web:15][web:18]

Assume that \(\bar{\varepsilon}(t)\) satisfies
\(0 \leq \bar{\varepsilon}(t) \leq \varepsilon_{\max}\). Then for any
\(c \in (0, 1/\varepsilon_{\max}]\) we obtain:
\[
  1 \leq D_s(t) \leq 2.
\]

\begin{lemma}[Uniform bounds on \(D_s(t)\)]
Suppose \(\bar{\varepsilon}(t) \in [0,\varepsilon_{\max}]\) for all \(t\)
and choose \(c = \frac{0.8}{\varepsilon_{\max}}\). Then the proxy
\(
  D_s(t) = 2 - 0.8 \bar{\varepsilon}(t)
\)
satisfies
\(
  1.2 \leq D_s(t) \leq 2
\),
with variation strictly controlled by the curvature range.
\end{lemma}

\begin{proof}
We have
\(
  D_s(t) \geq 2 - 0.8\varepsilon_{\max}
\)
and
\(
  D_s(t) \leq 2
\).
If \(\varepsilon_{\max} \leq 1\), then
\(
  D_s(t) \geq 2 - 0.8 = 1.2.
\)
\end{proof}

This elementary bound is sufficient to guarantee that the proxy
spectral dimension remains within the expected CDT range \([1,2]\) for
reasonable curvature scales, and that its oscillations or plateaus are
numerically controlled.

\section{Prime-Sector Dominance and Lifetimes}

Let \(t_n = n\,dt\) be discrete times and let \(p^\*(n)\) be the index of
the dominant sector at time \(t_n\), i.e.
\begin{equation}
  p^\*(n)
    := \arg\max_{p \in \mathcal{P}}
        \langle w \rangle_p(t_n),
  \qquad
  \langle w \rangle_p(t)
    = \frac{1}{|I_p|}\sum_{i \in I_p} w_{p,i}(t).
\end{equation}

For a fixed time window \(n \in \{n_0, \dots, N\}\), define the empirical
dominant-sector frequency
\begin{equation}
  \pi_p
    := \frac{1}{N - n_0 + 1}
       \bigl|\{n \in [n_0, N] : p^\*(n) = p\}\bigr|.
\end{equation}

\begin{proposition}[Prime-sector lifetime fraction]
For any finite set of primes \(\mathcal{P}\) and finite window
\([n_0,N]\), the vector \(\pi = (\pi_p)_{p \in \mathcal{P}}\) is a
probability vector:
\begin{equation}
  \pi_p \geq 0,
  \qquad
  \sum_{p \in \mathcal{P}} \pi_p = 1.
\end{equation}
Moreover, if the dynamics converges to a unique attracting set where a
single prime \(p_\infty\) dominates, then
\(
  \pi_p \to \delta_{p,p_\infty}
\)
as the window is shifted to late times.
\end{proposition}

\begin{proof}
By definition \(\pi_p\) counts the proportion of times in the discrete
set \(\{n_0,\dots,N\}\) for which prime \(p\) is dominant. Each time step
has exactly one dominant prime (assuming ties are broken consistently),
so \(\sum_p \pi_p = 1\). Nonnegativity is clear. If the dynamics
converges in the sense that there exists \(p_\infty\) such that
\(p^\*(n) = p_\infty\) for all sufficiently large \(n\), then any
sufficiently late window \([n_0,N]\) will have \(\pi_{p_\infty} \approx
1\), establishing the convergence.
\end{proof}

This provides a rigorous notion of \emph{sector lifetime fraction} and
connects the numerical experiments with the qualitative picture of prime
sector dominance and decoupling in the long-time limit.

\bigskip

\noindent
The results in this appendix formalize the operator-theoretic setting of
the prime-indexed noncommutative CDT model, provide explicit stability
and norm bounds for the feedback dynamics, and justify the numerical
choices used in the 2D toy simulator. They confirm that, under modest
assumptions on the Regge and NCG contributions, the multiplicity
evolution is well-posed and exhibits controlled behavior in both the
weight dynamics and the spectral-dimension proxy.[file:1][web:15][web:18]

\nocite{*}
\bibliographystyle{unsrt}  
\bibliography{references}
\end{document}